\documentclass[conference]{IEEEtran}

\usepackage{booktabs}
\usepackage{graphicx}
\usepackage{hyperref}
\usepackage{siunitx}

\title{Compact Molecular Descriptors for Honey Bee Toxicity Screening:\\
A Structure-Aware Comparison of Classical and Quantum Kernels}

% Double-blind review: restore author metadata only in the camera-ready copy.
\author{\IEEEauthorblockN{Anonymous Authors}}

\begin{document}
\maketitle

\begin{abstract}
This manuscript template describes a structure-aware study of acute pesticide
toxicity to \textit{Apis mellifera}. The study evaluates a compact ten-feature
molecular representation, controlled structural and partition-descriptor
ablations, and matched classical and simulated quantum kernels. Numerical
claims will be inserted only from versioned experiment bundles.
\end{abstract}

\begin{IEEEkeywords}
honey bee toxicity, molecular descriptors, pesticide screening, quantum
kernel, support vector machine, structure-aware validation
\end{IEEEkeywords}

\section{Introduction}
% Motivation: pollinator protection, limited species-specific data, and the
% value of compact/interpretable screening models.
% End with 3--4 concrete contributions, not a claim of quantum advantage.

\section{Related Work}
% ApisTox dataset and benchmark; molecular fingerprints and graph kernels;
% quantum kernels for compact molecular representations.

\section{Materials and Methods}
\subsection{Curated Molecular Domain and Endpoint}
% N=893; development N=712; historical holdout N=181; LD50 threshold;
% explain why this is a curated working domain rather than full ApisTox.

\subsection{X10 Molecular Representation}
% Define all ten descriptors, provenance, finite-value requirements, and the
% pre-specified X9/partition ablations.

\subsection{Structure-Aware Evaluation}
% Frozen Butina-cluster folds, nested inner model selection, no overlap,
% pooled OOF metrics, and status of the historical holdout.

\subsection{Classical Models}
% Logistic regression, RBF-SVC, random forest, grids, class weighting.

\subsection{Simulated Quantum Kernels}
% Exact statevectors, feature maps, scales, SVC with precomputed kernels,
% PSD checks, and matched molecule/feature/fold controls.

\subsection{Metrics and Statistical Analysis}
% Primary: AUROC. Secondary: AUPRC, balanced accuracy, MCC. Add paired
% uncertainty and kernel geometry diagnostics in phase 4.

\section{Results}
\subsection{X10 Classical Baseline}
% Generated from results/runs/<run_id>/summary.csv.

\subsection{Structural and Partition-Descriptor Ablations}

\subsection{Classical--Quantum Kernel Comparison}

\subsection{Structure-Disjoint Generalization}

\section{Discussion}
% Interpret signal, complementarity, model disagreement, practical limits,
% simulator-only scope, repeated holdout access, and descriptor availability.

\section{Conclusion}

\bibliographystyle{IEEEtran}
\bibliography{references}

\end{document}
